\documentclass[12pt]{article}
\usepackage[utf8]{inputenc}
\usepackage[T1]{fontenc}
\usepackage[margin=1in]{geometry}
\usepackage{setspace}
\usepackage{xcolor}
\usepackage[colorlinks=true, urlcolor=blue!60!black, linkcolor=blue!60!black]{hyperref}
\setlength{\parindent}{0em}
\setlength{\parskip}{0.6em}
\renewcommand{\baselinestretch}{1.0}
\begin{document}
\begin{center}
{\LARGE \textit{NLRB v. Jones \& Laughlin Steel Corp.}, 301 U.S. 1 (1937)\par}
\vspace{0.5em}
{\large Case Brief --- Two-Page Edition\par}
\end{center}

\section{Memory Jogger}
The 1937 turn: Congress may regulate labor relations in manufacturing that has a ``close and substantial relation to interstate commerce'' (\href{https://supreme.justia.com/cases/federal/us/301/1/}{301 U.S. 1}).

\section{Detailed Case Facts}
Jones \& Laughlin, the fourth-largest U.S. steel producer, ran an integrated enterprise of mines, ships, railroads, and plants across many states and shipped about 75\% of its steel out of Pennsylvania. It fired ten union leaders at its Aliquippa plant, and the NLRB, applying the Wagner Act, ordered reinstatement with back pay.

\section{Procedural History}
The Fifth Circuit refused to enforce the order as a regulation of local manufacturing. The Supreme Court granted certiorari and reversed.

\section{Judicial Votes}
Reversed, 5--4: Hughes wrote for Brandeis, Stone, Roberts, and Cardozo; McReynolds dissented for Van Devanter, Sutherland, and Butler.

\section{Holding}
The NLRA is valid as applied. Congress may regulate intrastate activities so closely and substantially related to interstate commerce that their control is essential or appropriate to protect it, and the Act's remedies satisfy due process and the Seventh Amendment.

\section{Analysis of Opinions}
Hughes read the Act to reach only practices ``affecting'' commerce, called the right to organize a ``fundamental right,'' and replaced formal categories with a practical test: ``It is the effect upon commerce, not the source of the injury, which is the criterion.'' Given the company's far-flung operations, a stoppage would be ``immediate'' and might be ``catastrophic.'' He distinguished \textit{Schechter} and \textit{Carter Coal}, but warned that federal power may not reach effects so indirect and remote as to erase the line between national and local, a caution later invoked in \textit{Lopez}. McReynolds insisted manufacturing is not commerce and that the effect of ten firings was ``indirect and remote in the highest degree.''

\section{Examples: Future Applications}
\textbf{Same-Side:} labor rules for a garment maker that buys and sells interstate are valid. \textbf{Opposite-Side:} a ban on guns near schools with no commerce hook exceeds the power, as in \textit{Lopez}. \textbf{Fence-Sitter:} a neighborhood bakery selling locally, reachable only through \textit{Wickard}-style aggregation.

\section{Critique}
Originalists such as Richard Epstein argue that ``commerce'' meant trade and that the effects test erased enumerated powers. Progressives praise its economic realism, while historians debate whether it was ``the switch in time that saved nine'' or organic doctrinal evolution.

\section{Key Quotations}
\begin{quote}
``Although activities may be intrastate in character when separately considered, if they have such a close and substantial relation to interstate commerce that their control is essential or appropriate to protect that commerce from burdens and obstructions, Congress cannot be denied the power to exercise that control.''
\end{quote}
\noindent --- Chief Justice Hughes, Opinion of the Court, 301 U.S. at 37.

\end{document}
